\documentclass[11pt,a4paper]{article}
\usepackage[utf8]{inputenc}
\usepackage[german]{babel}
\usepackage{amsmath,amssymb,amsthm}
\usepackage{geometry}
\geometry{margin=2.5cm}

\newtheorem{theorem}{Satz}[section]
\newtheorem{lemma}[theorem]{Lemma}
\newtheorem{definition}[theorem]{Definition}

\title{\textbf{Lokale Obstruktionen und globale diophantische Geometrie: Der Phasenübergang in der Defektschutz-Eigenschaft des Bamberger Modells für $p \equiv 7 \pmod 8$}}
\author{Thomas Hoffbauer \\ \small Bamberg, Deutschland}
\date{\today}

\begin{document}

\maketitle

\begin{abstract}
Wir untersuchen die Defektschutz-Eigenschaft (Low-Gap Protection Property) im Rahmen des Bamberger Modells für Primzahlen $p \equiv 7 \pmod 8$. Diese Eigenschaft postuliert, dass das Auftreten eines Defekts vom Typ $D$ mit Grad $g_D(p) = k$ das Vorhandensein eines Defekts vom Typ $O$ mit geringerem Grad ausschließt ($g_D(p) = k \implies g_O(p) \ge k$). Während für die ersten beiden Schichten ($k = 4$ und $k = 9$) ein vollständiger, rein lokal-arithmetischer Beweis ($L_4$ und $L_9$) durch Kombination archimedischer Produkt-Summen-Schranken und Restklassenobstruktionen ($\bmod 3, 7, 8, 16$) gelingt, zeigt sich bei $k = 36$ ($L_{36}$) ein fundamentaler methodischer Phasenübergang: Nach systematischer Filterung über Teiler, Paritäten und einen analytischen Minimalitätsfilter reduziert sich $L_{36}$ auf exakt sechs globale diophantische Restfälle. Für den ersten dieser Fälle zeigen wir die exakte Reduktion auf eine elliptische Quartik vom Geschlecht $1$ und ein gekoppeltes System verallgemeinerter Pell-Gleichungen sowie Lucas-Folgen.
\end{abstract}

\section{Einleitung und Grundgedanken}
Das Bamberger Modell untersucht die Struktur von Hurwitz-Koordinationen und Defekträumen über quadratischen Zahlkörpern und assoziierten Gitterstrukturen, insbesondere im Zusammenhang mit der Primzahlverteilung für $p \equiv 7 \pmod 8$. Ein zentrales Theorem betrifft die Stabilität kleiner Defekte, ausgedrückt durch Schutzsätze der Form:
\begin{equation}
g_D(p) = k \implies g_O(p) \ge k.
\end{equation}
Das System basiert auf Parametrisierungen der Form $2p = s^2 + \lambda$, wobei $\lambda$ durch die zugrundeliegenden Diagonalfamilien bestimmt wird.

\section{Die geschlossenen Sätze für $g_D = 4$ und $g_D = 9$ ($L_4$ und $L_9$)}
Für die Schichten $L_4$ und $L_9$ lässt sich zeigen, dass hypothetische Defekte $|\mathscr{D}| < k$ im $O$-Sektor unmöglich sind. Die Beweisarchitektur stützt sich auf:
\begin{enumerate}
    \item Archimedische Produkt-Summen-Abschätzungen für das Betragsverhältnis von $B$.
    \item Die Substitution $v = \frac{C+D}{2}$ und $z = \frac{C-D}{2}$.
    \item Lokale Restklassenobstruktionen modulo $3, 7, 8$ und $16$. Der Ast $v^2 - 2z^2 = -1$ erzwingt beispielsweise modulo 8 einen direkten Paritätswiderspruch.
\end{enumerate}

\section{Der Übergang zu $g_D = 36$ ($L_{36}$) und der Minimalitätsfilter}
Für $k = 36$ müssen die D36-Familien $\mathrm{D36}_{13}$ ($2p = s^2 + 13$) und $\mathrm{D36}_{37}$ ($2p = s^2 + 37$) getrennt betrachtet werden. Über die Quadratdifferenz $(t-s)(t+s) = \lambda - \mu$ schließt der Minimalitätsfilter störende kleinere Gaps aus und hebt die Startwerte an auf $x_{\min}^{(13)} = 81$ und $x_{\min}^{(37)} = 121$. Der Bereich $b > a$ kollabiert vollständig, wodurch $L_{36}$ auf sechs globale Restfälle reduziert wird.

\section{Analyse von Fall 1: Elliptische Quartik und Pell-Kompatibilität}
Für den ersten Restfall führt die rationale Parametrisierung $t = r/q$ auf die gemeinsame Quartik:
\begin{equation}
\mathcal{C}:\quad y^2 = 360x^4 + 60x^2 + 1.
\end{equation}
Über Pell-Nebenwirkungen und Einheiten in $\mathbb{Q}(\sqrt{10})$ reduziert sich das System auf die verallgemeinerte Pell-Gleichung $5W^2 - 2y^2 = 3$ und $W = 12x^2 + 1$, welche disjunkte Lucas-Folgen mit der Rekurrenz $T_{k+2} = 1442T_{k+1} - T_k$ erzeugen.

\section{Mathematische Anschlussgebiete und Historische Einordnung}
Die Arbeit verknüpft algebraische Zahlentheorie (quadratische Zahlkörper, Einheiten), diophantische Geometrie (elliptische Quartiken, Integralpunkte) und Rekurrenztheorie. Historisch reiht sie sich in die klassische diophantische Analysis ein und zeigt, dass vermeintlich statistische Primzahleigenschaften tiefen deterministischen Gesetzen der algebraischen Geometrie gehorchen.

\end{document}